\documentclass[twocolumn]{aastex631}
\begin{document}
\title{The reionization photon-budget ``crisis'' is a statement about $\xi_{\rm ion}$ and the star-formation-rate density: a systematic escape-fraction landscape}
\author{NebulaMind Lab (autonomous pipeline)}
\affiliation{NebulaMind Lab, \url{https://lab.nebulamind.net}}

\begin{abstract}
We ask whether star-forming galaxies alone can supply the ionizing photons required to reionize the Universe, and show that the answer is governed almost entirely by two adopted quantities---the ionizing-photon production efficiency $\xi_{\rm ion}$ and the cosmic star-formation-rate density (SFRD)---rather than by the escape fraction itself. Using the standard reionization-maintenance criterion, we solve for the Lyman-continuum escape fraction $f_{\rm esc}$ required of star-forming galaxies to balance recombinations across $5\le z\le 12$, propagating a grid of systematic assumptions ($\xi_{\rm ion}$, IGM clumping, LzLCS proxy calibration, and SFRD normalization) by Monte Carlo. The required $f_{\rm esc}$ rises steeply with redshift and crosses the indirect-proxy-inferred value ($\approx6\%$) at a ``crisis-onset'' redshift that ranges from $z\approx5$ to $z\approx8.75$ across defensible assumptions; it exceeds the physical ceiling $f_{\rm esc}=1$ at $z\approx8.5$ in the most pessimistic corner, but never within $5\le z\le12$ in the most optimistic. Under fiducial assumptions the budget closes to $z\approx6.5$ and becomes unphysical by $z\approx10.75$. We conclude that the apparent ``photon-budget crisis'' is not a robust observational result but a restatement of the adopted $\xi_{\rm ion}$ and SFRD priors: the same data admit both closure and a severe shortfall. This is a descriptive, automated study; it has not been validated by a human referee, uses no new survey data, and is not peer-reviewed.
\end{abstract}

\section{Introduction}
Reionization was largely complete by $z\approx5.3$, but which sources supplied the ionizing budget remains debated. A recurring claim is a ``photon-budget crisis'': the ionizing output of known star-forming galaxies appears insufficient to drive and maintain reionization, implying a shortfall that must be made up by additional sources or by unexpectedly high escape fractions \citep{Munoz2024,Davies2021}. Competing analyses argue that star-forming galaxies suffice once realistic ingredient values are adopted \citep{Duncan2015,Madau2017}. The comoving ionizing emissivity of galaxies is $\dot{n}_{\rm ion}=f_{\rm esc}\,\xi_{\rm ion}\,\rho_{\rm UV}$, where $\rho_{\rm UV}$ is the non-ionizing UV luminosity density (proportional to the SFRD), $\xi_{\rm ion}$ is the ionizing-photon production efficiency, and $f_{\rm esc}$ is the Lyman-continuum escape fraction \citep{Robertson2015}. Of these, $f_{\rm esc}$ is the least directly constrained: at $z>6$ the intergalactic medium is opaque to Lyman-continuum photons, so $f_{\rm esc}$ is inferred \emph{indirectly} from low-redshift-calibrated proxies (the O32 ratio, the UV continuum slope $\beta$; \citealt{Chisholm2022,Flury2022}), whose transportability to $z>6$ is uncertain. Because published ``required-$f_{\rm esc}$'' values are quoted under a single choice of $\xi_{\rm ion}$ and SFRD, the disagreement in the literature is, to first order, a disagreement in those two priors. Here we make that dependence explicit by computing the full systematic landscape of the required escape fraction, rather than a single point estimate, and by asking where---if anywhere---a genuine shortfall is robust.

\section{Method}
We adopt the reionization-maintenance criterion, in which the comoving ionizing emissivity balances recombinations, $\dot{n}_{\rm ion}=\langle n_{\rm H}\rangle/t_{\rm rec}$, with $t_{\rm rec}^{-1}=C\,\alpha_B\,n_{\rm H}(z)\,(1+Y/4X)$ \citep[case-B, $T=2\times10^4$\,K;][]{Robertson2015}. Solving $\dot{n}_{\rm ion}=f_{\rm esc}\,\xi_{\rm ion}\,\rho_{\rm SFR}/\kappa_{\rm UV}$ for the escape fraction gives the value $f_{\rm esc}(z)$ \emph{required} to close the budget. Inputs are literature-anchored and fixed---no survey catalog is queried: the SFRD is the \citet{MadauDickinson2014} analytic fit; the fiducial ionizing efficiency is $\log\xi_{\rm ion}=25.5\pm0.15$ \citep{Simmonds2024}; the clumping factor is $C\in[2,5]$; and $\kappa_{\rm UV}=1.15\times10^{-28}$. The proxy-inferred population escape fraction is taken from the LzLCS O32 and $\beta$ calibrations \citep{Chisholm2022,Flury2022}, with a median of $\approx6\%$ and $0.4$--$0.45$\,dex scatter. We evaluate a grid spanning $5\le z\le12$ in steps of $0.25$ crossed with eight systematic corners (a fiducial case; ``crisis'' and ``optimistic'' corners that push all systematics pessimistically or optimistically; and single-axis variants isolating $\xi_{\rm ion}$, clumping, the O32/$\beta$ calibration, and a JWST-enhanced SFRD). At each of 232 grid points a $4\times10^4$-sample Monte Carlo propagates the joint systematic uncertainty. As a non-circularity check we recompute the conclusion using the O32-only and $\beta$-only calibrations separately; the qualitative closure/shortfall verdict is unchanged by that swap at all redshifts. Finally we impose the physical bound $f_{\rm esc}\le1$: where the required value exceeds unity, star-forming galaxies cannot maintain reionization at \emph{any} escape fraction.

\section{Results}
Figure~\ref{fig:landscape} shows the required escape fraction as a function of redshift for all eight systematic corners. The required $f_{\rm esc}$ rises steeply with redshift---driven by the $(1+z)^3$ growth of the recombination rate and the falling SFRD---in every corner. Under fiducial assumptions it is $0.048$ at $z=6$, $0.21$ at $z=8$, and $0.685$ at $z=10$, crossing the proxy-inferred value ($\approx0.06$) at a crisis-onset redshift $z\approx6.5$ and the physical ceiling $f_{\rm esc}=1$ at $z\approx10.75$. Across the systematic corners the crisis-onset redshift spans $z\approx5.0$ (the ``crisis'' corner: low $\xi_{\rm ion}$, high clumping, $\beta$ calibration, no SFRD enhancement) to $z\approx8.75$ (the ``optimistic'' corner: high $\xi_{\rm ion}$, low clumping, O32 calibration, JWST-enhanced SFRD). The unphysical boundary ($f_{\rm esc}>1$) ranges from $z\approx8.5$ in the crisis corner to \emph{never} (within $5\le z\le12$) in the optimistic and JWST-high-SFRD corners. The dominant levers are $\xi_{\rm ion}$ and the SFRD normalization; the choice of proxy calibration (O32 versus $\beta$) shifts the inferred value only modestly and does not change the qualitative picture.

\begin{figure}
\centering
\includegraphics[width=\columnwidth]{landscape.png}
\caption{Escape fraction required of star-forming galaxies to close the reionization ionizing budget versus redshift, for eight systematic corners (grey band: full envelope). The dash-dotted line is the indirect-proxy-inferred value; the dotted line is the physical ceiling $f_{\rm esc}=1$. The redshift at which each corner crosses the inferred value (crisis onset) and unity (galaxies-cannot) varies by several redshift units across defensible assumptions.\label{fig:landscape}}
\end{figure}

\section{Discussion}
The two positions in the literature---that star-forming galaxies fall short \citep{Munoz2024,Davies2021} and that they suffice \citep{Duncan2015}---are \emph{both} recovered within our systematic envelope: they correspond to pessimistic and optimistic choices of $\xi_{\rm ion}$ and the SFRD, respectively. The data therefore do not robustly favour either. What would break the degeneracy is not a better escape-fraction proxy but tighter constraints on the two dominant ingredients: direct $\xi_{\rm ion}$ measurements from hydrogen recombination lines at high redshift, independent SFRD normalizations, and lower-redshift Lyman-continuum anchors to test proxy transportability. Practically, a stated ``photon-budget crisis'' should always be accompanied by the assumed $\xi_{\rm ion}$ and SFRD, since the crisis-onset redshift moves by $\sim3.5$ redshift units as those assumptions vary over their plausible ranges. We stress two structural limitations that bound the interpretation. First, the escape fraction is inferred \emph{entirely} from indirect, low-redshift-calibrated proxies; if $f_{\rm esc}$ genuinely rises toward high redshift, as some models predict, the crisis-onset would move to higher $z$ than any corner here captures, and only direct high-redshift Lyman-continuum or $\xi_{\rm ion}$ measurements can settle it. Second, the maintenance criterion assumes ionization equilibrium rather than integrating the full reionization history, which sets the normalization of the required emissivity. Neither limitation alters the central result---that the closure verdict is dominated by $\xi_{\rm ion}$ and the SFRD, not by $f_{\rm esc}$---but both caution against reading any single crisis-onset redshift as definitive.

\section{Caveats}
This is an automated, descriptive study and has not been validated by a human referee. The proxy-inferred escape fraction is held redshift-independent; a genuinely rising $f_{\rm esc}(z)$, as some models predict, would delay the crisis onset. We adopt a single SFRD parameterization; the JWST high-redshift SFRD is itself contested and is bracketed here only by the enhancement corner. The maintenance criterion assumes ionization equilibrium rather than integrating the full reionization history. All escape fractions are inferred from indirect, low-redshift-calibrated proxies; no direct $z>6$ Lyman-continuum detection is used. Required values exceeding unity are reported as a diagnostic of the shortfall magnitude and are physically inadmissible as escape fractions.

\begin{thebibliography}{}
\bibitem[Chisholm et al.(2022)]{Chisholm2022} Chisholm, J., et al. 2022, MNRAS, 517, 5104
\bibitem[Davies et al.(2021)]{Davies2021} Davies, F.~B., et al. 2021, ApJ, 918, L35
\bibitem[Duncan \& Conselice(2015)]{Duncan2015} Duncan, K., \& Conselice, C.~J. 2015, MNRAS, 451, 2030
\bibitem[Flury et al.(2022)]{Flury2022} Flury, S.~R., et al. 2022, ApJ, 930, 126
\bibitem[Madau(2017)]{Madau2017} Madau, P. 2017, ApJ, 851, 50
\bibitem[Madau \& Dickinson(2014)]{MadauDickinson2014} Madau, P., \& Dickinson, M. 2014, ARA\&A, 52, 415
\bibitem[Mu\~noz et al.(2024)]{Munoz2024} Mu\~noz, J.~B., et al. 2024, MNRAS, 535, L37
\bibitem[Robertson et al.(2015)]{Robertson2015} Robertson, B.~E., et al. 2015, ApJ, 802, L19
\bibitem[Simmonds et al.(2024)]{Simmonds2024} Simmonds, C., et al. 2024, MNRAS, 527, 6139
\end{thebibliography}
\end{document}
